\subsection{Training and Inference Strategy}

The overall optimization of the proposed zero-shot fault diagnosis framework is directed by a composite objective function that integrates generative quality, semantic alignment, and structural reconstruction. Following the individual loss components defined in Sections 3.6 and 3.7, the total training loss $\mathcal{L}_{total}$ is formulated as:
\begin{equation}
\mathcal{L}_{total} = \mathcal{L}_{diffusion} + \lambda_1 \mathcal{L}_{AC} + \lambda_2 \mathcal{L}_{recon}
\end{equation}
where $\mathcal{L}_{diffusion}$ denotes the standard variational lower bound for the reverse diffusion process, $\mathcal{L}_{AC}$ is the Attribute Consistency Loss ensuring semantic anchoring, and $\mathcal{L}_{recon}$ represents the reconstruction error for the latent feature extractor. The hyperparameters $\lambda_1$ and $\lambda_2$ serve as scaling factors to balance the contribution of the semantic regularizer and the reconstruction fidelity, respectively.

\subsubsection{Optimization and Hyperparameters}
To achieve stable convergence amidst the high-dimensional vibration data, we employ the AdamW optimizer, which incorporates decoupled weight decay to prevent overfitting in the generative path. The learning rate follows a cosine annealing schedule, gradually decaying from an initial value of $2 \times 10^{-4}$ to a minimum of $1 \times 10^{-6}$ over the course of training. Initial experiments showed that without cosine annealing, the model tended to converge to sub-optimal local minima for the TEP dataset. Key hyperparameters for the diffusion process include the number of diffusion steps $T=1000$ and a linear variance schedule $\beta_t$. We explored the use of a cosine variance schedule but found that the linear schedule performed better for periodic mechanical signals. The embedding dimension $d$ for the vibration features is set to 512 to match the CLIP semantic space, while $\lambda_1$ and $\lambda_2$ are empirically tuned to 0.1 and 0.05, respectively.

\subsubsection{Two-Phase Training Procedure}
The training process is executed in two distinct phases to ensure the robustness of the synthesized features. In Phase I: Feature Path Pretraining, the latent feature extractor and the attribute projection network $f(\cdot)$ are pretrained on the seen classes to establish a reliable mapping between the raw vibration signals and the CLIP-aligned feature space. This phase typically lasts for 100 epochs with a batch size of 64. In Phase II: Joint Fine-tuning, the diffusion-based generator $G$ is integrated, and the entire pipeline is optimized end-to-end using $\mathcal{L}_{total}$ for another 300 epochs. This phased approach prevents the generator from collapsing during the early stages of training when the feature extractor has not yet stabilized. During Phase II, we use a smaller learning rate ($5 \times 10^{-5}$) for the pre-trained components to preserve the established feature space while allowing the diffusion path to align with it.

\subsubsection{Inference for Unseen Fault Diagnosis}
During the inference stage, the model performs zero-shot prediction for previously unseen fault categories. Given the semantic embedding $\mathbf{s}_{unseen}$ of a target unseen class, the reverse diffusion process is sampled to generate a set of synthetic vibration features $\hat{\mathbf{x}}_{unseen}$ for that class. Usually, we generate $N=500$ synthetic samples per unseen class to capture the expected intra-class variation. These features are then used to train a simple softmax classifier or a distance-based prototype matcher. For a test sample $\mathbf{x}_{test}$, the trained feature extractor maps it to the latent space, and the final classification is determined by calculating the cosine similarity between the test feature and the centroids of the synthesized unseen class distributions. This pipeline enables the model to generalize its diagnostic capabilities to novel fault conditions without requiring any labeled temporal data from those classes during training. Our experimental results show that this sampling-based inference approach is significantly more accurate than using a single point-prototype for the unseen classes.
